\section{Related Work}

Our research builds upon and extends work in several domains: computer vision for crop disease detection, explainable AI techniques, and agricultural decision support systems. This section summarizes relevant prior work in these areas and identifies the gaps that our research addresses.

\subsection{Computer Vision for Crop Disease Detection}

Deep learning approaches have revolutionized crop disease detection in recent years. The seminal work by \cite{mohanty2016using} demonstrated the potential of CNNs for plant disease classification, achieving 99.35\% accuracy on the PlantVillage dataset using transfer learning with architectures like AlexNet and GoogLeNet. This research established a baseline for subsequent work but noted significant limitations in real-field conditions.

Focusing specifically on paddy diseases, \cite{lu2017identification} proposed a deep CNN framework that achieved 95.48\% accuracy across ten rice disease categories. \cite{rahman2020deep} further improved accuracy to 97.5\% by optimizing hyperparameters and data augmentation techniques. More recently, \cite{singh2023comparative} evaluated various CNN architectures for paddy disease classification, finding that ensemble approaches could reach 98.2\% accuracy on the Paddy Doctor dataset.

Object detection frameworks have also been applied to localize and identify multiple diseases within a single image. \cite{fuentes2017robust} compared Faster R-CNN, R-FCN, and SSD for plant disease detection, while \cite{liu2020yolo} adapted YOLOv3 for real-time disease detection in field conditions, achieving 91.0\% mAP with 18 FPS processing speed.

Despite these impressive accuracy metrics, a common limitation across these studies is their emphasis on prediction performance without addressing model interpretability. Most implementations treat the deep learning model as a black box, providing little insight into why a particular disease was detected or which visual features contributed to the classification decision.

\subsection{Explainable AI Techniques}

Explainable AI (XAI) has emerged as a critical research area aimed at making AI systems more transparent and interpretable. Approaches to explainability can be broadly categorized as inherently interpretable models and post-hoc explanation techniques for black-box models.

Among post-hoc visual explanation techniques for CNNs, Gradient-weighted Class Activation Mapping (Grad-CAM) \cite{selvaraju2017grad} has gained widespread adoption due to its ability to produce class-specific visual explanations without requiring architectural changes. Other notable approaches include LIME \cite{ribeiro2016should}, which approximates local model behavior with interpretable surrogates, and SHAP \cite{lundberg2017unified}, which assigns feature importance values based on game-theoretic principles.

For uncertainty quantification in deep learning, \cite{gal2016dropout} demonstrated that Monte Carlo Dropout provides a practical approximation of Bayesian inference. This technique enables neural networks to express uncertainty in their predictions—a critical capability for high-stakes applications like disease diagnosis. \cite{kendall2017uncertainties} further distinguished between aleatoric uncertainty (inherent data noise) and epistemic uncertainty (model uncertainty), showing how both can be estimated and communicated.

In agricultural contexts, \cite{obermeyer2019dissecting} explored how explanation techniques could be adapted for crop disease detection, but focused primarily on technical accuracy rather than agricultural utility. \cite{phillips2020adaptive} proposed a domain-specific explanation framework for plant stress detection but did not integrate uncertainty information or connect explanations to treatment recommendations.

\subsection{Agricultural Decision Support Systems}

Agricultural Decision Support Systems (DSS) aim to assist farmers in making informed management decisions based on data analysis. Traditional DSS like DSSAT \cite{jones2003dssat} and APSIM \cite{keating2003overview} have focused on simulation-based approaches that model crop growth and response to management interventions. These systems typically offer transparent reasoning but lack the pattern recognition capabilities of modern deep learning.

More recent AI-enhanced DSS have incorporated machine learning to provide data-driven recommendations. \cite{zhang2019crop} developed a recommendation system for optimizing fertilizer application based on yield prediction models. \cite{liakos2018machine} reviewed machine learning applications in agriculture, noting the trade-off between predictive power and interpretability across various applications.

Mobile applications for crop disease diagnosis have gained popularity, with platforms like Plantix and PlantVillage offering AI-powered identification. However, \cite{wright2022adoption} found that farmers' trust in these applications was limited by the lack of explanations and contextual sensitivity. \cite{german2021co} emphasized the importance of co-designing agricultural AI systems with farmers to ensure recommendations align with local knowledge and practices.

\subsection{Research Gaps and Our Contributions}

From our analysis of related work, we identify several critical gaps that our research addresses:

\begin{enumerate}
    \item \textbf{Integration Gap}: Most prior work treats disease detection, explanation generation, and recommendation systems as separate problems rather than as integrated components of a unified decision support framework.
    
    \item \textbf{Contextual Gap}: Existing explanation techniques rarely incorporate domain-specific agricultural knowledge, limiting their relevance and utility for farmers and extension workers.
    
    \item \textbf{Uncertainty Gap}: Few disease detection systems explicitly communicate uncertainty alongside predictions, despite the inherent variability in field conditions and image quality.
    
    \item \textbf{Actionability Gap}: There is limited research on connecting visual explanations to practical treatment and management recommendations in a context-aware manner.
    
    \item \textbf{Evaluation Gap}: Most explainable AI research in agriculture evaluates technical metrics (explanation faithfulness, localization) without assessing impact on user trust and decision quality.
\end{enumerate}

Our ExplainableCropAI framework addresses these gaps through a holistic approach that combines advanced visual explanations with uncertainty quantification and domain-specific knowledge to generate actionable recommendations. Unlike previous work that focuses primarily on technical explanation quality, we emphasize the practical utility of explanations in agricultural decision-making contexts.
